%% LyX 2.4.4 created this file.  For more info, see https://www.lyx.org/.
%% Do not edit unless you really know what you are doing.
\documentclass[english,t]{beamer}
\usepackage{mathptmx}
\usepackage[T1]{fontenc}
\usepackage[latin9]{inputenc}
\usepackage{amssymb}
\usepackage{cancel}

\makeatletter
%%%%%%%%%%%%%%%%%%%%%%%%%%%%%% Textclass specific LaTeX commands.
% this default might be overridden by plain title style
\newcommand\makebeamertitle{\frame{\maketitle}}%
% (ERT) argument for the TOC
\AtBeginDocument{%
  \let\origtableofcontents=\tableofcontents
  \def\tableofcontents{\@ifnextchar[{\origtableofcontents}{\gobbletableofcontents}}
  \def\gobbletableofcontents#1{\origtableofcontents}
}

%%%%%%%%%%%%%%%%%%%%%%%%%%%%%% User specified LaTeX commands.
\usetheme{Warsaw}
% or ...

\setbeamercovered{transparent}
% or whatever (possibly just delete it)
\setbeamercovered{invisible} %makes overlays invisible


%gets rid of navigation symbols
\setbeamertemplate{navigation symbols}{}

\usepackage{multicol}

\usepackage{tikz}
\usetikzlibrary{calc}
\usetikzlibrary{plotmarks}
\usepackage{pgfplots}
\pgfplotsset{compat=newest}

\usepackage{calc}
\usepackage[nomessages]{fp}

\usepackage{advdate}    % Advancing/saving dates
\usepackage{datetime}   % Dates formatting
\usepackage{datenumber} % Counters for dates

%\usepackage{pgfpages}
%\pgfpagesuselayout{4 on 1}[a4paper, landscape, border shrink=7mm]
%\pgfpageslogicalpageoptions{1}{border code=\pgfusepath{stroke}}
%\pgfpageslogicalpageoptions{2}{border code=\pgfusepath{stroke}}
%\pgfpageslogicalpageoptions{3}{border code=\pgfusepath{stroke}}
%\pgfpageslogicalpageoptions{4}{border code=\pgfusepath{stroke}}
%%%Don't forget to add 'handout'

\makeatother

\usepackage{babel}
\begin{document}
\newcommand{\xmax}{2000}
\newcommand{\xmin}{0}
\newcommand{\xstep}{0} %must be xmin+n
\newcommand{\ymax}{500}
\newcommand{\ymin}{0}
\newcommand{\ystep}{0} %must be ymin+n

\newcounter{countEx} \setcounter{countEx}{0}
\newcounter{countProb} \setcounter{countProb}{0}
\resetcounteronoverlays{countEx} \resetcounteronoverlays{countProb}

\newcommand{\ex}[3]{
	\addtocounter{countEx}{1}
	\only<#1-#2>{\begin{exampleblock} {Example \chNum.\secNum.\arabic{countEx}.} #3	\end{exampleblock}}
}

\newcommand{\prob}[4]{
	\addtocounter{countProb}{1}
	\only<#1-#2>{\begin{exampleblock} {Problem  \chNum.\secNum.\arabic{countProb}.\,#3} #4	\end{exampleblock}}
}

\newcommand{\yline}[4]{
	(#2 - #4)/(#1 - #3)*(x - #1) + #2
}

\beamerdefaultoverlayspecification{<+->}

%%%%Chapter and Section numbers%%%%%
\newcommand{\chNum}{2}
\newcommand{\secNum}{2}
\newcommand{\secTitle}{Solving Systems of Linear Equations by the Gauss-Jordan Method}

\title[Section \chNum.\secNum: \secTitle]{Section \chNum.\secNum: \secTitle}

\subtitle{Finite Mathematics~~MTH 132~~Spring 2014}

\author{Dr.\,James Ashe}

\titlegraphic{\includegraphics[height=1.5cm]{/home/jim/JamesAshe/JCSU/images/jcsu_bull}}

\institute{Johnson C. Smith University}

\date{{}}

\makebeamertitle 
\begin{frame}{Solving Linear Systems by the Gauss-Jordan method}

\begin{columns}


\column{5.5cm}
\begin{itemize}
\item <1->[]$\begin{array}{rrrrr}
-2x & + & y & = & -8\\
2x & + & y & = & 4
\end{array}$
\item <3->[]$\begin{array}{rrrrr}
-2 & \,\,\,\,\,\,\,\,\, & 1 & \,\,\,\,\,\,\,\,\, & -8\\
2 & \,\,\,\,\,\, & 1 & \,\,\,\, & 4
\end{array}$
\item <4->[]$\begin{bmatrix}\begin{array}{rr|r}
-2 & 1 & -8\\
2 & 1 & 4
\end{array}\end{bmatrix}$
\end{itemize}

\column{5.5cm}

\vspace{-.525cm}
\begin{block}<2->{step 1}
Write the system in abbreviated form, called an \emph{augmented
matrix}.
\end{block}
\end{columns}
\end{frame}

\begin{frame}{}

\begin{columns}


\column{5.5cm}
\begin{itemize}
\item <1->[]$\begin{bmatrix}\begin{array}{rr|r}
-2 & 1 & -8\\
2 & 1 & 4
\end{array}\end{bmatrix}$
\end{itemize}

\column{5.5cm}

\vspace{-.525cm}
\begin{block}<2->{step 2}
Transform the matrix so that each column (except the last) has no
more than one $0$.
\end{block}

\begin{block}<3->{step 3}
Divide so the leading enrty for each row is $1$.
\end{block}

\begin{block}<4->{step 4}
A row of $0$'s means no solution. Otherwise the last column gives
the solution.
\end{block}
\end{columns}
\end{frame}

\begin{frame}{}

\ex{1}{}{Solve the system $\begin{array}{rrrrr} -3x & + & y & = & 4\\ 2x & - & 2y & = & -4 \end{array}$}
\end{frame}

\begin{frame}{}

\ex{1}{}{Use the Gauss-Jordan method to solve the system.}
\begin{itemize}
\item <1->[]$\begin{array}{rrrr}
x &  & +z & =3\\
 & y & -z & =1\\
x & +y & +z & =4
\end{array}$
\item <2->[]$\begin{bmatrix}\begin{array}{rrr|r}
1 & 0 & 1 & 3\\
0 & 1 & -1 & 1\\
1 & 1 & 1 & 4
\end{array}\end{bmatrix}$
\end{itemize}
\end{frame}

\begin{frame}{}

\prob{1}{}{Use the Gauss-Jordan method to solve the system.}{\hfill $\begin{array}{rrrrr} x & + & y & = & 5\\ 3x & + & 2y & = & 12 \end{array}$ \hfill}
\begin{itemize}
\item <2->[]Answer: $(2,3)$
\end{itemize}
\end{frame}

\end{document}
